%
% Chapter 3 — System Requirements
%
\chapter{System Requirements} \label{cap:requirements}

This chapter presents the functional and non-functional requirements of the Diário LX platform, derived from the project specification and validated with the editorial team.

\section{Functional Requirements} \label{sec:req_functional}

\subsection{Content Management} \label{sec:req_content}

The system must provide a full content management lifecycle supporting the following content types: \textbf{articles}, \textbf{podcasts}, \textbf{videos}, and \textbf{photo essays}. Each content item must include the following attributes:

\begin{itemize}
  \item Title
  \item Body text (supporting inline subtitles, highlights, and captioned images)
  \item Lead paragraph (\textit{entrada})
  \item Author(s) and photo author
  \item Creation and publication date
  \item Featured image
  \item Category and subcategory
  \item Tags
\end{itemize}

Content must progress through a defined workflow:

\begin{enumerate}
  \item \textbf{Draft} — created and edited by a Contributor.
  \item \textbf{Pending Approval} — submitted for editorial review.
  \item \textbf{Published} — approved and publicly visible.
  \item \textbf{Rejected} — returned to the author with revision comments.
\end{enumerate}

It must be possible to archive and permanently delete content. Media file uploads are limited to 2\,GB. It must be possible to mark content as featured for display on the homepage.

\subsection{Categories, Subcategories, and Tags} \label{sec:req_categories}

Content must be organised using categories and subcategories. The main editorial sections are:

\begin{itemize}
  \item Lisboa, Cidade Aberta
  \item A Fundo
  \item Secções (Mundo, Política, Sociedade, Economia, Cultura e Media, Bem-Viver, Desporto)
  \item Especiais
  \item Fotografia
  \item Podcasts
  \item Vídeos
\end{itemize}

Each category must support a name, an associated colour, and an automatically generated listing page. Tags must also have dedicated listing pages. The system must allow the creation, editing, and deletion of categories and tags.

\subsection{User Management and Permissions} \label{sec:req_users}

The system defines three user roles:

\begin{description}
  \item[Administrator] Full system control: manages users, system settings, approves or rejects content and new accounts, and has access to all backoffice areas.
  \item[Editor] Can approve or reject submitted content, manage contributor accounts, and deactivate users.
  \item[Contributor] Can create and submit content for approval. Cannot publish directly.
\end{description}

Each user profile includes: name, email, profile photograph, active/inactive status, professional function (journalist or photographer), and role.

\subsection{Backoffice} \label{sec:req_backoffice}

The backoffice must provide:

\begin{itemize}
  \item Account creation for Editors and Contributors.
  \item Login and logout for all profiles.
  \item Listing and search of all articles and users.
  \item Profile editing (own profile for all; any profile for Administrators).
  \item Comment moderation (approve, reject, hide) by Administrators and Editors.
\end{itemize}

\subsection{Public Website} \label{sec:req_frontend}

\subsubsection{Homepage}
The homepage must display:
\begin{itemize}
  \item 4 featured articles (defined in the backoffice).
  \item The 4 most recent articles.
  \item Navigation to main editorial sections.
  \item Links to institutional pages: Team, Editorial Statute, and Code of Ethics.
\end{itemize}

\subsubsection{Other Pages}
\begin{itemize}
  \item \textbf{Category page}: one featured article (defined in the backoffice) and remaining articles in reverse chronological order.
  \item \textbf{Tag page}: list of associated articles.
  \item \textbf{Article page}: full article view, reader comments, and 4 related articles.
  \item \textbf{Team page}.
  \item \textbf{Podcast page}: dedicated layout distinct from other category pages.
\end{itemize}

\subsection{Search and Navigation} \label{sec:req_search}

\begin{itemize}
  \item Full-text keyword search across articles.
  \item Human-readable, slugified URLs (e.g. \texttt{/politica/orcamento-estado-2026}).
\end{itemize}

\subsection{Compatibility} \label{sec:req_compat}

The system must be compatible with modern browsers and fully functional on mobile devices (responsive design for mobile, tablet, and desktop viewports).

\section{Non-Functional Requirements} \label{sec:req_nonfunctional}

\begin{itemize}
  \item \textbf{Security}: Authentication via token-based mechanism; role-based access control enforced at the API level.
  \item \textbf{Performance}: Pages must load within acceptable response times under expected editorial team and reader load.
  \item \textbf{Maintainability}: Codebase organised to allow independent evolution of frontend and backend.
  \item \textbf{Media storage}: Support for images, audio (podcasts), and video files up to 2\,GB per upload.
  \item \textbf{Scalability}: Architecture must support growth in both content volume and concurrent users.
\end{itemize}

\section{Optional Requirements} \label{sec:req_optional}

The following features are considered optional and may be addressed in future iterations:

\begin{itemize}
  \item Scheduled publication of content.
  \item Advanced search filters (by date, category, author, tag) in both backoffice and public site.
  \item User search filters in the backoffice (by role, function, status).
  \item Reader account registration and commenting.
  \item Audiodescription support for visually impaired users.
  \item Multilingual support (Portuguese and English).
\end{itemize}
